% GENERATED FILE. Edit canonical lessons, then run npm run generate:books.
% canonical-source-hash: fnv1a64:319de9f494c08fda
% canonical-lessons: JA-C184-shinamono, JA-C184-retsu, JA-C184-torikaeru, JA-C184-zenbu, JA-C184-choudo

\chapter{Goods, Queues and Exchanges}
\label{ch:ja-goods-queues-and-exchanges}
\hlchaptermodality{184}

\begin{chapteropening}
\textbf{By the end of this chapter:} \emph{I can talk about goods, queues and totals.}

Complete the last lesson of chapter 184: I can talk about goods, queues and totals.
\end{chapteropening}

\section[\texorpdfstring{\ja{しなもの} (shinamono)}{しなもの (shinamono)}]{\ja{しなもの} (shinamono) — goods; article, item}

\label{lesson:JA-C184-shinamono}



\begin{quote}
Before the new one: say the Japanese for to choose, then the Japanese for to line up.
\end{quote}

\subsection*{You'll want to know: \ja{しなもの}}
\textbf{\ja{しなもの}} — \emph{shinamono} — \textquotedblleft{}goods; article, item\textquotedblright{}.

\begin{quote}
\textbf{\ja{この} \ja{みせの} \ja{しなものは} \ja{やすいです。}} — \emph{kono mise no shinamono wa yasui desu} — \textquotedblleft{}The goods in this shop are cheap.\textquotedblright{}
\end{quote}

\begin{cousinweb}
\textbf{\ja{しなもの}} (\emph{shinamono}) is native Japanese, built from \emph{shina}, \textquotedblleft{}article, goods\textquotedblright{}, and \emph{mono}, \textquotedblleft{}thing\textquotedblright{}, both inherited rather than borrowed (JMdict, KANJIDIC2).
\end{cousinweb}

\begin{grammarlens}[title={What you've built}]
One more word for everyday life at A2.
\end{grammarlens}

\subsection*{Your turn}
\begin{itemize}
\raggedright
  \item \emph{Say it:} \emph{shinamono}
  \item \emph{Say it:} \emph{shinamono}, once more
  \item \emph{Say it:} \emph{kono mise no shinamono wa yasui desu} — The goods in this shop are cheap.
  \item \emph{Recall:} say the Japanese for to choose, then the Japanese for to line up, then say \emph{shinamono} again
\end{itemize}

\begin{tcolorbox}[breakable,colback=teal!4,colframe=teal!35!black,title={Before you move on}]
What does \ja{しなもの} mean? (\textquotedblleft{}Goods; article, item\textquotedblright{}.) Say it once more.
\end{tcolorbox}
\section[\texorpdfstring{\ja{れつ} (retsu)}{れつ (retsu)}]{\ja{れつ} (retsu) — queue; line, row}

\label{lesson:JA-C184-retsu}



\begin{quote}
Before the new one: say the Japanese for to line up, then the Japanese for goods.
\end{quote}

\subsection*{You'll want to know: \ja{れつ}}
\textbf{\ja{れつ}} — \emph{retsu} — \textquotedblleft{}queue; line, row\textquotedblright{}.

\begin{quote}
\textbf{\ja{みせの} \ja{まえに} \ja{ながい} \ja{れつが} \ja{あります。}} — \emph{mise no mae ni nagai retsu ga arimasu} — \textquotedblleft{}There is a long queue in front of the shop.\textquotedblright{}
\end{quote}

\begin{cousinweb}
\textbf{\ja{れつ}} (\emph{retsu}) is Sino-Japanese (kango): one Chinese character, \textquotedblleft{}row, file\textquotedblright{}, read with its on'yomi, a reading borrowed from Chinese (Middle Chinese \emph{ljet}) (KANJIDIC2; Middle Chinese in Baxter's transcription).
\end{cousinweb}

\begin{grammarlens}[title={What you've built}]
One more word for everyday life at A2.
\end{grammarlens}

\subsection*{Your turn}
\begin{itemize}
\raggedright
  \item \emph{Say it:} \emph{retsu}
  \item \emph{Say it:} \emph{retsu}, once more
  \item \emph{Say it:} \emph{mise no mae ni nagai retsu ga arimasu} — There is a long queue in front of the shop.
  \item \emph{Recall:} say the Japanese for to line up, then the Japanese for goods, then say \emph{retsu} again
\end{itemize}

\begin{tcolorbox}[breakable,colback=teal!4,colframe=teal!35!black,title={Before you move on}]
What does \ja{れつ} mean? (\textquotedblleft{}Queue; line, row\textquotedblright{}.) Say it once more.
\end{tcolorbox}
\section[\texorpdfstring{\ja{とりかえる} (torikaeru)}{とりかえる (torikaeru)}]{\ja{とりかえる} (torikaeru) — to exchange; to swap, to replace}

\label{lesson:JA-C184-torikaeru}



\begin{quote}
Before the new one: say the Japanese for goods, then the Japanese for queue.
\end{quote}

\subsection*{You'll want to know: \ja{とりかえる}}
\textbf{\ja{とりかえる}} — \emph{torikaeru} — \textquotedblleft{}to exchange; to swap, to replace\textquotedblright{}.

\begin{quote}
\textbf{\ja{てんいんが} \ja{くつを} \ja{とりかえます。}} — \emph{ten'in ga kutsu o torikaemasu} — \textquotedblleft{}The shop assistant exchanges the shoes.\textquotedblright{}
\end{quote}

Polite \textbf{\ja{とりかえます}}, \emph{torikaemasu}; an ichidan verb.

\begin{cousinweb}
\textbf{\ja{とりかえる}} (\emph{torikaeru}) is native Japanese, built from \emph{tori}, the stem of \emph{toru}, \textquotedblleft{}to take\textquotedblright{}, and \emph{kaeru}, \textquotedblleft{}to exchange\textquotedblright{}, both inherited rather than borrowed (JMdict, KANJIDIC2).
\end{cousinweb}

\begin{grammarlens}[title={What you've built}]
One more word for everyday life at A2.
\end{grammarlens}

\subsection*{Your turn}
\begin{itemize}
\raggedright
  \item \emph{Say it:} \emph{torikaeru}
  \item \emph{Say it:} \emph{torikaeru}, once more
  \item \emph{Say it:} \emph{ten'in ga kutsu o torikaemasu} — The shop assistant exchanges the shoes.
  \item \emph{Recall:} say the Japanese for goods, then the Japanese for queue, then say \emph{torikaeru} again
\end{itemize}

\begin{tcolorbox}[breakable,colback=teal!4,colframe=teal!35!black,title={Before you move on}]
What does \ja{とりかえる} mean? (\textquotedblleft{}To exchange; to swap, to replace\textquotedblright{}.) Say it once more.
\end{tcolorbox}
\section[\texorpdfstring{\ja{ぜんぶ} (zenbu)}{ぜんぶ (zenbu)}]{\ja{ぜんぶ} (zenbu) — all of it; everything, in all}

\label{lesson:JA-C184-zenbu}



\begin{quote}
Before the new one: say the Japanese for queue, then the Japanese for to exchange.
\end{quote}

\subsection*{You'll want to know: \ja{ぜんぶ}}
\textbf{\ja{ぜんぶ}} — \emph{zenbu} — \textquotedblleft{}all of it; everything, in all\textquotedblright{}.

\begin{quote}
\textbf{\ja{ぜんぶで} \ja{せんえんです。}} — \emph{zenbu de sen'en desu} — \textquotedblleft{}It comes to a thousand yen in all.\textquotedblright{}
\end{quote}

\begin{cousinweb}
\textbf{\ja{ぜんぶ}} (\emph{zenbu}) is Sino-Japanese (kango): two Chinese characters read with their on'yomi, readings borrowed from Chinese. They are \emph{zen}, \textquotedblleft{}whole\textquotedblright{} (Middle Chinese \emph{dzjwen}) and \emph{bu}, \textquotedblleft{}part, section\textquotedblright{} (\emph{buX}) (KANJIDIC2; Middle Chinese in Baxter's transcription).
\end{cousinweb}

\begin{grammarlens}[title={What you've built}]
One more word for everyday life at A2.
\end{grammarlens}

\subsection*{Your turn}
\begin{itemize}
\raggedright
  \item \emph{Say it:} \emph{zenbu}
  \item \emph{Say it:} \emph{zenbu}, once more
  \item \emph{Say it:} \emph{zenbu de sen'en desu} — It comes to a thousand yen in all.
  \item \emph{Recall:} say the Japanese for queue, then the Japanese for to exchange, then say \emph{zenbu} again
\end{itemize}

\begin{tcolorbox}[breakable,colback=teal!4,colframe=teal!35!black,title={Before you move on}]
What does \ja{ぜんぶ} mean? (\textquotedblleft{}All of it; everything, in all\textquotedblright{}.) Say it once more.
\end{tcolorbox}
\section[\texorpdfstring{\ja{ちょうど} (chōdo)}{ちょうど (chōdo)}]{\ja{ちょうど} (chōdo) — exactly; precisely, right}

\label{lesson:JA-C184-choudo}



\begin{quote}
Before the new one: say the Japanese for to exchange, then the Japanese for all of it.
\end{quote}

\subsection*{You'll want to know: \ja{ちょうど}}
\textbf{\ja{ちょうど}} — \emph{chōdo} — \textquotedblleft{}exactly; precisely, right\textquotedblright{}.

\begin{quote}
\textbf{\ja{ちょうど} \ja{せんえんです。}} — \emph{chōdo sen'en desu} — \textquotedblleft{}It is exactly a thousand yen.\textquotedblright{}
\end{quote}

\begin{cousinweb}
\textbf{\ja{ちょうど}} (\emph{chōdo}) is Sino-Japanese (kango): two Chinese characters read with their on'yomi, readings borrowed from Chinese. They are \emph{chō} (Middle Chinese \emph{teng}) and \emph{do}, \textquotedblleft{}degree\textquotedblright{} (\emph{duH}); JMdict marks it as usually written in kana (JMdict, KANJIDIC2).
\end{cousinweb}

\begin{grammarlens}[title={What you've built}]
One more word for everyday life at A2.
\end{grammarlens}

\subsection*{Your turn}
\begin{itemize}
\raggedright
  \item \emph{Say it:} \emph{chōdo}
  \item \emph{Say it:} \emph{chōdo}, once more
  \item \emph{Say it:} \emph{chōdo sen'en desu} — It is exactly a thousand yen.
  \item \emph{Recall:} say the Japanese for to exchange, then the Japanese for all of it, then say \emph{chōdo} again
\end{itemize}

\begin{tcolorbox}[breakable,colback=teal!4,colframe=teal!35!black,title={Before you move on}]
What does \ja{ちょうど} mean? (\textquotedblleft{}Exactly; precisely, right\textquotedblright{}.) Say it once more.
\end{tcolorbox}
